\documentclass{amsart}
%\usepackage[backend=bibtex8,style=numeric,giveninits=true]{biblatex}
%\addbibresource{formschub.bib}
\usepackage{algorithm}
\usepackage{algpseudocode}
\usepackage{hyperref}
\usepackage{fancyvrb}
\usepackage{amsmath}
\usepackage{amsfonts}
\usepackage{amsthm}
\usepackage{amssymb}
\usepackage{array}
\usepackage{cases}
\usepackage{caption}
\usepackage{xcolor}
\usepackage{tikz}
\usepackage{tikz-cd}
\usepackage{mathtools}
\usetikzlibrary{calc}
\usepackage{centernot}
\usepackage{mathtools}
\usepackage{graphicx}
\usepackage[margin=1in]{geometry}
\usepackage[makeroom]{cancel}
\usepackage{stmaryrd}
\usepackage{stackengine}
\usepackage{mathrsfs}
\usepackage{enumitem}

% Concatenation product on \dcoma: a square cup with a + inside (visually `\sqcupplus`).
\makeatletter
\newcommand{\sqcupplus@inner}[2]{%
  \vcenter{\hbox{\ooalign{%
    \hfil$\m@th#1\sqcup$\hfil\cr
    \hfil\raisebox{0.15ex}{$\m@th#1\scriptscriptstyle+$}\hfil\cr
  }}}%
}
\newcommand{\sqcupplus}{\mathbin{\mathpalette\sqcupplus@inner\relax}}
\makeatother

%\usepackage{unicode-math}
%\usepackage{nicematrix}

%\usepackage{witharrows}

\newtheorem{lemma}{Lemma}[subsection]
\newtheorem{corollary}[lemma]{Corollary}

\newtheorem{proposition}[lemma]{Proposition} 
\newtheorem{proposition2}{Proposition}[section]
\newtheorem{theorem}{Theorem}[section]
\newtheorem{conjecture}[theorem]{Conjecture} 
\newtheorem{observation}{Observation}[section]
\newtheorem{question}{Question}

%\newtheorem{conj}[thm]{Conjecture} 


%%% 
%%% The following gives definition type environments (which only differ
%%% from theorem type invironmants in the choices of fonts).  The
%%% numbering is still tied to the theorem counter.
%%% 

\theoremstyle{definition}
%\newtheorem{definition}{Definition}[section]
%\newtheorem{example}{Example}[section]
\newtheorem{definition}[lemma]{Definition}
\newtheorem{example}[lemma]{Example}
%\newtheorem{algorithm}{Algorithm}
\newtheorem*{note}{Note}
\newtheorem*{method}{Method}
%%
%% Again more of these can be added by uncommenting and editing the
%% following. 
%%

%\newtheorem{note}[thm]{Note}


%%% 
%%% The following gives remark type environments (which only differ
%%% from theorem type invironmants in the choices of fonts).  The
%%% numbering is still tied to the theorem counter.
%%% 


\theoremstyle{remark}

\newtheorem*{remark}{Remark}


%%%
%%% The following, if uncommented, numbers equations within sections.
%%% 

%\numberwithin{equation}{section}
\numberwithin{equation}{subsection}


%%%
%%% The following show how to make definition (also called macros or
%%% abbreviations).  For example to use get a bold face R for use to
%%% name the real numbers the command is $\mathbf{R}.  To save typing we
%%% can abbreviate as

\newcommand{\R}{\mathbf{R}}  % The real numbers.

%%
%% The comment after the definition is not required, but if you are
%% working with someone they will likely thank you for explaining your
%% definition.  
%%
%% Now add you own definitions:
%%

%%%
%%% Mathematical operators (things like sin and cos which are used as
%%% functions and have slightly different spacing when typeset than
%%% variables are defined as follows:
%%%

\DeclareMathOperator{\sch}{\mathfrak{S}}
\newcommand{\shup}{{\uparrow}}
\newcommand{\downvar}[1]{\stackrel{#1}{\searrow}}
\newcommand{\wof}[1]{w_{#1}}
\newcommand{\rcdd}{\mathfrak{F}}
\newcommand{\wtt}{\mathtt{wt}}
\newcommand{\QY}{\ensuremath{\mathbb{QY}}}
\newcommand{\tom}[1]{\xrightarrow{#1}}
\newcommand{\Tom}[1]{\xRightarrow{#1}}
\newcommand{\mot}[1]{\xleftarrow{#1}}
\newcommand{\grass}{\mathcal{G}}
\newcommand{\plac}{\mathrm{Plac}}

\newcommand{\cperm}[1]{\llbracket #1\rrbracket}
\newcommand{\cpcf}[5]{d^{#1,#2}_{#3}\left(#4\mid #5\right)}
% simplified product notation
\newcommand{\smpfu}{\mathsf{\Pi}}
\newcommand{\smpr}[3]{\ensuremath{}\smpfu\left( #1 \mid #2_{[#3]}\right)}
%\newcommand{\smpre}[2]{\ensuremath{}\smpfu\left( #1 \mid #2\right)}
\newcommand{\smpre}[2]{\ensuremath{}\prod_{b\in #2}(#1-b)}

\newcommand{\rtt}{\mathtt{rt}}
\newcommand{\trm}{\mathtt{trim}}
\newcommand{\zeromap}{\mathcal{Z}}
\newcommand{\hr}{\mathtt{ht}}
\newcommand{\clp}{\mathtt{clip}}
\newcommand{\BRC}{\mathcal{BRC}}
\newcommand{\RC}{\mathcal{RC}}
\newcommand{\lzeromap}{\mathscr{Z}}
\newcommand{\lmap}{\mathscr{L}}
\newcommand{\ltb}{\mathrm{little}}
\newcommand{\slide}{\mathfrak{F}}
%\newcommand{\shifthom}{\operatornamewithlimits{\raisebox{1ex}{$\stackrel{\blacktriangleright}{\mathbf{\_}}$}}\limits}

%\newcommand{\shifthom}{\operatornamewithlimits{\raisebox{1ex}{$\blacktriangleright$}}\limits}

\newcommand{\shifthom}{\operatornamewithlimits{\triangleright}\limits}


%\newcommand{\shiftvby}[1]{\ensuremath{\resizebox{\width}{\heightoff}{$\shifthom_{#1}$}}}

\newcommand{\shiftvby}[1]{\shifthom_{#1}}

\newcommand*\circled[1]{%
	\tikz[baseline=(char.base)]{
		\node[shape=circle,draw,inner sep=0.1pt] (char) {#1};
	}%
}
\newcommand\mycircled[1]{\makebox[0pt]{\circled{#1}}}
\newcommand{\tb}{\textbullet}
\newcommand{\xr}[1]{\ensuremath{{}\xrightarrow{[#1]}{}}}
\newcommand{\rd}[1]{\ensuremath{\mathcolor{red}{\mathbf{#1}}}}
\newcommand{\mc}[1]{\ensuremath{\mycircled{#1}}}
\newcommand{\desc}{\mathrm{Desc}}%
\newcommand{\SSYT}{\mathrm{SSYT}}
%\newcommand{\schv}[2]{\operatornamewithlimits{\sch_{\mathit{#2}}}\limits_{\negphantom{{}_{\mathit{#2}}}\negvphantom{{}_{\mathit{#2}}}#1}}

%\newcommand{\schv}[2]{\setstackgap{S}{0pt}\setstackgap{L}{0pt}\stackMath\stackunder{\sch_{#2}}{\negphantom{\scriptstyle{#2}}\!{\scriptscriptstyle{#1}}}}

\newcommand{\schv}[2]{\shiftvby{#1}\!\sch_{#2}}

%\newcommand{\shiftvby}[1]{\makebox[0pt]{$\triangleright$}#1}



%\newglossary[slg]{symbolslist}{syi}{syg}{List of symbols}
%\makeglossaries

%\begin{itemize}
%	\item Elements of $S_\infty$ $c{[}p,q{]}$, $d{[}p,q{]}$, $r{[}p,q{]}$, homomorphism $\sigma_p:S_\infty\to S_\infty$: Definition \ref{definition:specelem}	

%	\item  Relations between elements of $S\infty$ $\tom{k},\Tom{k}$: Definition \ref{definition:pierisymb}

%	\item Code, dual code, $\code(w)$, $\code^*(w)$, dominant permutations, the dominant approximation $\dom(v)$: Definition \ref{definition:codestuff}

%	\item Function/sets of permutations related to pulling out variables from Schubert polynomials $\varphi_{i,n}$, $\mathcal{D}_i(v)$, $Q_i(v',v)$: Definition \ref{definition:polysequence}

%	\item Flattening function $\phi_i:S_\infty\to S_\infty$, domination relation between permutations: Definition \ref{definition:domination}

%	\item Special polynomials $H_p(x;y)$, $E_p(x;y)$, $\smpr{x_1}{y}{p}$: Definition \ref{definition:hp}

%	\item Polynomial variable omission $x^{(i)}$, subscript substitution $y_A$, index shift function $\shiftvby{i}$: Definition \ref{definition:pv}
%\end{itemize}

\newcommand{\dsch}{\ensuremath{\Xi}}
\newcommand{\coma}{\mathcal{A}}
\newcommand{\dcoma}{\mathcal{D}}



%%
%% This is the end of the preamble.
%% 
\title{A Littlewood--Richardson rule for forest polynomials via the Schubert bialgebra}
\author{Matthew J. Samuel}

%\newcommand{\tomm}[2]{\xrightarrow{#1}_{#2}}
%\counterwithout{equation}{section}
%\setcounter{section}{-1}
\newcommand{\shuff}[2]{\mathcal{S}\mathcal{H}_{#1}^{#2}}
\newcommand{\dom}{\mathfrak{d}}
\newcommand{\ldom}{\dom^*}
\newcommand{\code}{\mathfrak{c}}
\newcommand{\ducode}{\overline{\mathfrak{c}}}
\newcommand{\forest}{\mathfrak{P}}
\newcommand{\fcode}{\mathfrak{c}_\forest}
\newcommand{\finv}{P_{\forest}}

\newcommand{\maxd}{\mathfrak{m}}
%\newcommand{\ddv}[1]{{{}_{{\scriptscriptstyle (#1)}}\partial}}
\newcommand{\ddv}[1]{\delta}
%\newglossary[slg]{symbolslist}{syi}{syg}{Symbolslist}
%\GlsXtrLoadResources
%[src={symbols},sort=use,type=main,
%	group=specelem]
%\glsaddstoragekey{group}{}{\grouplabel}
%\glsxtrsetgrouptitle{specelem}{}





%\newglossary[slg]{symbol}{sot}{stn}{Symbols}
%\makeglossaries

\begin{document}

\begin{abstract}
\end{abstract}

	\maketitle


\section{Introduction}


\section{Set-valued inversions tableaux}
\begin{definition}

Let $w$ be a permutation such that $\maxd(w)\leq n$. Suppose $T:I(w)\to \mathbf{2}^{[n]}\setminus \{\varnothing\}$ is a function mapping a root to a nonempty subset of $[n]$. Define relations $\leq$ and $<$ on these sets by declaring that $A\leq B$ if $\max(A)\leq \min(B)$ and $A<B$ if $\max(A)<\min(B)$. 
We say that $T$ is a \emph{set-valued inversions tableau} if the following conditions are satisfied:
\begin{itemize}
  \item If $i < j < k$ and $(i,j),(j,k),(i,k)\in I(w)$, then either $T(i,j) \leq T(i,k) \leq T(j,k)$ or $T(j,k) \leq T(i,k) \leq T(i,j)$.
  \item If $i>0$ and $j<k$ are such that $(i,j),(i,k)\in I(w)$, then $T(i,j) < T(i,k)$.
  \item $T(i,i+1)\leq \{i\}$ for all $i$.

\end{itemize}
For $(i,j)\in I(w)$, define $\lambda_T(i,j)$ as follows. Suppose $T(i,j) = r$. Let $(i,k)$ be the root such that $T(i,k)=r$ and $k$ is minimal. Then $\lambda_T(i,j) = i + j + 1 - k$.
\end{definition}

\begin{theorem}
Suppose $w$ is a permutation. Then we have the following formula for the double Grothendieck polynomial of $w$:
$$\mathfrak{G}_w^{(\beta)}(x;y) = \sum_{T} \prod_{(i,j)\in I(w)} \beta^{|T(i,j)| - 1} \prod_{r\in T(i,j)} x_r \oplus y_{\lambda_T(i,j) - r}$$
where the sum is over all set-valued inversions tableaux of $w$. 
\end{theorem}

\begin{definition}
An \emph{staircase word graph} is simply a subset of $\mathbb{N}^2$. We impose the same ordering on these as RC graphs. That is, we say that $(a,b)\leq (c,d)$ if $a<c$ or $a=c$ and $b > d$. We define $\mathbf{word}(G)$ to be the word obtained by reading the elements $i + j - 1$ in the linear order given by $\leq$. We then say that $\wof{G}$ is the Demazure product of $\mathbf{word}(G)$.

For a set-valued inversions tableau $T$, we can define a staircase word graph $G$ as follows. If $w$ is the identity, then $G$ is the empty set. Otherwise, let $n$ be the maximum value that is contained in $T(i,j)$ for some $(i,j)\in I(w)$. Choose $i$ and $j$ such that $T(i,j)=n$ and $i + j$ is as small as possible. Necessarily, $j=i+1$, or in other words $i$ is a right descent of $w$. If $|T(i,i+1)| > 1$, remove $n$ from $T(i,i+1)$ to obtain a new function $T':I(w)\to \mathbf{2}^{[n]}\setminus \{\varnothing\}$. If $|T(i,i+1)|=1$, then let $T':I(ws_i)\to \mathbf{2}^{[n]}\setminus \{\varnothing\}$ be the function defined by
$$T'(a,b) = T(s_i(a),s_i(b))$$
for all $(a,b)\in I(ws_i)$.

It is straightforward to check that $T'$ is a set-valued inversions tableau, either for $w$ if $|T(i,i+1)| > 1$ or for $ws_n$ if $|T(i,i+1)| = 1$. Let $G'=\Phi(T')$. Then let $\Phi(T)=G'\cup \{(n,i+1-n)\}$. 
\end{definition}

\begin{proposition}
$\Phi$ is a bijection between set-valued inversions tableaux of $w$ with maximum value $n$ and staircase word graphs $G$ such that $\wof{G} = w$ and $a\leq n$ for all $(a,b)\in G$.
\end{proposition}
\begin{proof}
  If $|T(i,j)|=1$ for all $(i,j)\in I(w)$, then $T$ is just an ordinary inversion tableau and $\Phi(T)$ is just the usual RC graph associated to $T$. In this case, the result is \cite{axelrodfreed2025inversionstableaux}. Otherwise, we prove this by induction on the maximum size of $|T(i,j)|$. Suppose this maximum size is $m>1$ and we have proved the result for all set-valued inversions tableaux with maximum size less than $m$. We may assume without loss of generality that $m$ is the maximum value that is contained in $T(i,j)$ for some $(i,j)\in I(w)$ by removing elements below. In that case, $m\in T(i,i+1)$ for some $i$. Let $T'$ be the set-valued inversions tableau obtained by removing $m$ from $T(i,i+1)$ if $|T(i,i+1)| > 1$ or by applying the transformation described in the definition of $\Phi$ if $|T(i,i+1)|=1$. By the inductive hypothesis, $\Phi(T')$ is a staircase word graph such that $\wof{\Phi(T')} = w$ or $\wof{\Phi(T')} = ws_i$ and $a\leq n$ for all $(a,b)\in \Phi(T')$. In either case, we have $\wof{\Phi(T)} = w$. It is also clear that $a\leq n$ for all $(a,b)\in \Phi(T)$. This proves that $\Phi(T)$ is a staircase word graph with the desired properties.

  Let $G$ be a staircase word graph. Define $\Psi(G)$ to be empty if $G$ is empty. Otherwise, let $(i,j)$ be the maximum element of $G$ and let $G' = G\setminus \{(i,j)\}$. Let $w'=\wof{G'}$. Let $d=i + j - 1$. If $\ell(w's_d) > \ell(w')$, then let
$$\Psi(G)(a,b) = \begin{cases}
\Psi(G')(s_d(a),s_d(b)) & \text{if } (a,b)\in I(w)\setminus\{(i,j)\}\\
\{d\} & \text{if } (a,b) = (i,j)
\end{cases}$$
If $\ell(w's_d) < \ell(w')$, then let
$$\Psi(G)(a,b) = \begin{cases}
\Psi(G')(a,b) & \text{if } (a,b)\in I(w)\setminus\{(i,j)\}\\
\Psi(G')(a,b)\cup \{d\} & \text{if } (a,b) = (i,j)
\end{cases}$$
\end{proof}

\begin{definition}
Define the isobaric divided difference operator
$$\pi^i = \partial^i(1+\beta x_{i+1})$$
Note the Leibniz formula
$$\pi^i(fg) = (\pi^i+\beta s_i)(f)g + f\pi^i(g)$$
\end{definition}


\cite{samuelmolev}
\bibliographystyle{acm}
\bibliography{schubert_bialgebra}


\end{document}
